% Imporditud: tools/import_eksperimendid.py
% Allikas: Eesti füüsikaolümpiaadi eksperimendi ülesannete
%   kogu (Rael Kalda, 2023), ülesanne Ü34, lahendus L34.
\setAuthor{EFO žürii}
\setRound{piirkonnavoor}
\setYear{1997}
\setNumber{G E1}
\setDifficulty{2}
\setTopic{vedelike mehaanika}
\setVahendid{mõõtesilinder, vesi, puitklots, kivi, mõõtejoonlaud}
\setPunktid{12}
\setAllikas{Eksperimendi kogu Ü34, lahendus lk 47}
\setLahendusseis{toores}

\prob{Puitklotsi ja kivi tihedus}
Leida puitklotsi ja kivi tihedused. Hinnata mõõteviga.

\hint

\solu
Teooria:

Keha ujub, kui FU = mg , kus FU on üleslükke jõud, m on keha mass ja g on vaba langemise kiirendus (1 p). FU = $\rho$V g , kus $\rho$ on vee tihedus, V on välja tõrjutud vee ruumala (1 p). Ujuva keha mass võrdub välja tõrjutud vee massiga m = $\rho$V (1 p).

Suhtelise vea valem: E(x) = $\Delta$x/x (1 p);

Summa (vahe) arvutamise viga: E ($x_1$ $\pm$ $x_2$) = ($\Delta$$x_1$ + $\Delta$$x_2$) / ($x_1$ $\pm$ $x_2$) (1 p);

Korrutise (jagatise) arvutamise viga: E ($x_1$/$x_2$) = E ($x_1$ $\cdot$ $x_2$) = E ($x_1$) + E ($x_2$) (1

p).

Mõõtmised: Mõõdame puitklotsi mõõtmed (1 p). Mõõdame, kui suure vee ruumala VP V tõrjub välja ujuv puitklots (1 p). Mõõdame, kui suure vee ruumala VK tõrjub välja kivi (1 p). Mõõdame, kui suure vee ruumala VKP tõrjub välja kivi, ujudes puitklotsil (1 p).

Arvutused: Puitklotsi ruumala on VP = a $\cdot$ b $\cdot$ c, kus a, b ja c on puitklotsi mõõtmed. Puitklotsi mass on mP = VP V $\cdot\rho$V , kus $\rho$V on vee tihedus. Puitklotsi tihedus on $\rho$P = mP /VP . Kivi mass on mK = (VKP $-$ VP V ) $\cdot \rho$V . Kivi tihedus on $\rho$K = mK /VK .

Mõõtevea arvutamine: Mõõtmiste suhtelised vead (1 p): Puuklotsi mõõtmete mõõtmise suhtelised vead: E(a) = $\Delta$a/a, E(b) = $\Delta$b/b, E(c) = $\Delta$c/c. Välja tõrjutud vee ruumala mõõtmise suhteline viga: E(V ) = 2$\Delta$V /V , kus V on vastav mõõdetud välja tõrjutud vee ruumala ja $\Delta$V on antud mõõtesilindri mõõtetäpsus.

Arvutuste suhtelised vead (1 p): Puitklotsi ruumala arvutamise suhteline viga: E (VP ) = E(a) + E(b) + E(c). E (mP ) = E (VP ) --- puitklotsi massi suhteline viga on võrdne puitklotsi ruumala suhtelise veaga. Puitklotsi tiheduse arvutamise suhteline viga on võrdne puitklotsi ruumala suhtelise vea ja puitklotsi poolt välja tõrjutud vee ruumala suhtelise vea summaga: E ($\rho$P ) = E (VP ) + E (VP V ). Kivi massi arvutamise suhteline viga: E (mK ) = 2$\Delta$V / (VKP $-$ VP V ), kus $\Delta$V on mõõtesilindri mõõtetäpsus. Kivi tiheduse arvutamise suhteline viga on võrdne kivi massi arvutamise suhtelise vea ja kivi ruumala mõõtmise suhtelise vea summaga: E ($\rho$K ) = E (mK ) + E (VK ).
\probend
